\chapter{The Interpolation Problem}
\label{chap:interpolation-problem}

% Closed question: What can be said about the gap between recorded events
% without inventing structure not in the ledger?

\begin{chapteressay}

The Lean elaborator performs a computation that has no obvious counterpart
in human reasoning. When a user types
\begin{verbatim}
def x := f y
\end{verbatim}
without specifying the type, the elaborator must figure out what type
\texttt{x} should have. It does this by introducing a metavariable: a
named placeholder for the unknown type. The metavariable is not zero. It
is not unspecified. It is a specific token, recorded in the elaborator's
state, with a list of constraints attached. As the elaborator processes
more of the program, the constraints accumulate. At some point, the
constraints uniquely determine the metavariable's value, and the
elaborator commits.

The remarkable property is that this can take a long time. A definition
that contains many implicit arguments, complicated typeclass resolutions,
or recursive references may introduce dozens of metavariables. Each one
must be carried through the elaboration, constrained by the surrounding
context, and resolved at the right moment. If the elaborator resolves
too early, it commits to a wrong value and produces a confusing error
later. If it resolves too late, the elaboration may stall waiting for
constraints that never arrive. The elaborator's algorithm is a careful
choreography of constraint gathering and metavariable resolution.

This choreography is the algorithmic embodiment of the chapter's
question. The compiler is between two recorded events: the user's
source code (the body of the definition) and the final elaborated
term (the type-checked, fully resolved value that will be added to
the namespace). Between those two events, the elaborator carries a
soup of partial information: types being inferred, instances being
resolved, universes being unified. The carriage is admissible only if
the partial information is held honestly --- as constrained
placeholders, not as committed but possibly wrong values.

This chapter's claim is that any algorithm that operates on incomplete
information must use this structure. The algorithm cannot avoid having
gaps; the gaps are inherent in operating before all input has been
seen. The question is how the algorithm represents the gaps. The
chapter develops three concrete examples, each from a different domain,
showing how the metavariable structure appears.

The first example is cubic spline fitting. Ten temperature readings
have been taken at ten times during the day. The question ``what was
the temperature at 9:15 AM, between the 8 AM and 10 AM readings?'' has
no direct answer; the instrument did not run at 9:15. The cubic spline
algorithm fills the gap with a uniquely determined smooth curve through
the data. The algorithm's resolution of the gap is the algorithm's
output. The interior slopes of the spline are metavariables; they are
resolved by a linear system derived from the smoothness conditions.

The second example is the Lean elaborator itself. The user types a
definition with a hole, and the elaborator infers the missing type
from the body and the context. The hole is a metavariable. The
inference is the resolution. The elaborator carries the metavariable
through the elaboration, gathering constraints, and resolves it when
the constraints are sufficient. The same algorithmic structure as
the cubic spline, but in the domain of types rather than the domain
of temperatures.

The third example is the boundary of what these algorithms can do.
Some metavariables cannot be resolved. The bits of Chaitin's
constant $\Omega$ are well-defined mathematically, constrained by
the halting status of each program, but no algorithm can compute them.
The metavariable can be carried by the elaborator, but cannot be
evaluated. This is the halting boundary, encountered as the
interpolation limit.

The chapter's second metaphor is a legal contract. A contract is a
\texttt{REPRESENTABLE} legal instrument: well-formed, signed
according to specified rules, witnessed by a notary. A contract may
have blanks: a date to be filled in, signatures to be obtained,
initial fields awaiting their owner. The contract carries the blanks
as metavariables. Each blank has a type (a date, a signature, an
initial) and admissibility conditions (within a time window,
witnessed, executed by a specific party). The contract is
admissible in the unsigned state; it is also admissible in the
signed state. The transition is the resolution of the metavariables.

The lawyer drafting the contract does not silently fill in blanks
with default values. She does not pretend the contract is complete.
She manages the metavariables: she identifies what each blank
requires, she gathers the constraints, she sequences the resolutions.
The contract may close successfully (all blanks resolved, signatures
witnessed) or may halt unresolved (offer expires). Either outcome is
admissible. The record carries the resolution status honestly.

The chapter's closed question is:

\begin{center}
\emph{What can be said about the gap between recorded events without
inventing structure not in the ledger?}
\end{center}

The chapter develops the answer through five sections. Section 1
introduces cubic spline fitting as the canonical algorithmic
interpolation. Section 2 examines the Lean elaborator's metavariable
machinery as the chapter's running illustration. Section 3 introduces
the halting boundary as the limit case: a metavariable that no finite
algorithm can resolve. Section 4 covers noisy interpolation: when the
data itself is uncertain, the metavariable is resolved to an
estimate with named uncertainty. Section 5 examines
\texttt{ChaitinsNumberSequence} in Lean: the cleanest example of a
type that exists but admits no constructible value, the
representation of the unresolvable metavariable.

The structure throughout is the same. The algorithm has incomplete
information. It carries the incomplete information as a metavariable.
It accumulates constraints from the surrounding context. When the
constraints suffice, it resolves; when they don't, it carries the
metavariable forward, marked as unresolved. The honest report of the
resolution status is part of the algorithm's output.

This pattern is the chapter's contribution to the book's larger
argument. The compiler is doing this. The universe, if it is running
admissible algorithms, must be doing this. The selection rule that
picks one physical history from many admissible ones is a
metavariable resolution problem at the deepest level: the
constraint to be satisfied is the exact Kolmogorov complexity of the
history, and the metavariable to be resolved is which of the
infinitely many admissible histories has the minimum complexity.
Chapter 6 takes up that argument. The earlier chapters build the
vocabulary.

\end{chapteressay}

\section{Cubic Spline Fitting}
\label{se:cubic-spline-fitting}

Ten temperature readings have been taken at ten times during the day: 6 AM,
8 AM, 10 AM, noon, 2 PM, 4 PM, 6 PM, 8 PM, 10 PM, midnight. The readings
are integers in degrees Celsius. The question is: what was the temperature
at 9:15 AM? No reading was taken at that time, but the value must be
something. The instrument was not running at 9:15; the ledger has no entry.
The task is to fill in a value that is admissible given the surrounding
readings, without inventing data the instrument did not produce.

The Python library \texttt{scipy.interpolate} provides
\texttt{CubicSpline} for exactly this task. Given the ten (time,
temperature) pairs, the constructor builds a piecewise cubic function that
passes through every pair and is twice continuously differentiable
($C^2$) at the interior knots. Evaluating the spline at 9:15 AM produces a
specific number; the spline's continuity properties guarantee that the
number is consistent with the surrounding data.

The construction proceeds as follows. On each of the nine intervals
between adjacent readings, the spline is a cubic polynomial $a_i + b_i
(t - t_i) + c_i (t - t_i)^2 + d_i (t - t_i)^3$ for $i = 0, 1, \ldots, 8$.
There are nine intervals and four coefficients per interval, giving thirty-
six unknowns. The constraints are: the spline passes through both
endpoints of each interval (two equations per interval, eighteen total);
first derivatives match at each interior knot (eight equations); second
derivatives match at each interior knot (eight equations). Total
constraints: thirty-four. Two more constraints are needed; the standard
choice is the natural spline boundary, $f''(t_0) = f''(t_9) = 0$.

The thirty-six unknowns and thirty-six equations form a linear system.
The system is tridiagonal in the slopes at the interior knots, which
means it can be solved in $O(n)$ time by Gaussian elimination. The
solution is unique: there is exactly one $C^2$ cubic spline through the
data with natural boundary conditions. The spline is the algorithm's
output.

What was the temperature at 9:15 AM? The spline evaluated at 9:15
produces, say, twenty-one degrees Celsius. This is the algorithm's
answer. The answer is not a measurement; the instrument did not take a
reading at 9:15. The answer is an interpolation: the unique smooth value
consistent with the surrounding measurements under the cubic-spline
admissibility rule.

The distinction between measurement and interpolation is the chapter's
subject. The ten readings are measurements: the instrument produced
them, the ledger admits them, they are facts in the record. The
spline's value at 9:15 is an inference: the algorithm produces it from
the surrounding facts under a specified continuity assumption. The
inference is admissible only because the continuity assumption is
stated explicitly. If the assumption changes (linear interpolation,
piecewise constant, quadratic, fifth-order polynomial), the inferred
value changes. The interpolation depends on the assumption; the
measurement does not.

The cubic spline is the canonical algorithmic answer because it is the
smoothest interpolation that does not invent excess structure. Higher-
order polynomials (a single ninth-degree polynomial through all ten
points) would oscillate wildly between the data points; the Runge
phenomenon is well-documented. Lower-order interpolations (linear,
piecewise quadratic) lack second-derivative continuity, which produces
visible kinks at the knots. Cubic splines achieve $C^2$ smoothness
with the minimum polynomial degree, and the natural boundary condition
minimizes the integral of the squared second derivative ---the bending
energy of the curve.

This is the curve-fitting analog of the chapter's question. The gap
between two recorded events (the temperature at 8 AM and the
temperature at 10 AM) is filled by the unique $C^2$ piecewise cubic
that minimizes bending energy. The fill is structured; it is not
arbitrary. The structure is named (cubic spline with natural boundary).
The structure is admissible because it does not introduce features
(oscillations, kinks, discontinuities) that the data does not
require.

In Lean, the analog is the \texttt{Spline} type from \texttt{Episode6.lean}
in the Lean codebase. The \texttt{Spline} structure carries a list of
(time, value, slope) triples; the slope at each interior knot is the
metavariable that the spline-fitting algorithm resolves. The natural
boundary condition fixes the slope at the endpoints. The interior
slopes are determined by the tridiagonal solve. The \texttt{Spline.le}
relation is the partial order on splines: one spline refines another
when it has the same or more knot points with the same or more
specific slopes.

The cubic spline fitting is the chapter's first concrete example of an
algorithm that carries a metavariable. Before the linear system is
solved, the interior slopes are unknown. They are not arbitrary: the
constraints determine them. They are not measured: the instrument did
not produce them. They are inferred: the algorithm produces them from
the constraints. Carrying the metavariable through the algorithm,
without resolving it prematurely, is what permits the eventual unique
solution.

\section{Carrying Metavariables in the Lean Elaborator}
\label{se:carrying-metavariables}

Open a Lean file and type:
\begin{verbatim}
def foo : ?$\alpha$ := 3
\end{verbatim}
The expression \texttt{?$\alpha$} is a placeholder for the type. The user
has not specified what type \texttt{foo} should have; they have asked Lean
to figure it out. The Lean elaborator processes the definition as follows:
it sees the body \texttt{3}, recognizes it as a numeric literal, infers
that the type must be a numeric type, and resolves \texttt{?$\alpha$} to
\texttt{Nat} (the default numeric type unless overridden). The definition
is admitted with type \texttt{Nat}.

The placeholder \texttt{?$\alpha$} is a metavariable. It is the
elaborator's record that some information is needed but not yet
available. During elaboration, the metavariable is carried as a typed
hole: the elaborator knows that \texttt{?$\alpha$} is a type (of sort
\texttt{Type 0}), but does not know which type. As elaboration proceeds,
constraints accumulate: the body imposes that \texttt{?$\alpha$} admits
the literal \texttt{3}; the type-class resolution imposes that
\texttt{?$\alpha$} has a \texttt{OfNat 3} instance; the default-instance
rule imposes that, absent other constraints, \texttt{?$\alpha$} is
\texttt{Nat}. The constraints, taken together, uniquely determine
\texttt{?$\alpha$ = Nat}.

Turn on the elaboration trace with \texttt{set\_option
trace.Meta.isDefEq true} and re-run. The trace prints every unification
step: every time the elaborator asks ``is this metavariable equal to this
other term?'', the trace records the question and the answer. The trace
for the above definition shows two or three unification events, each one
constraining \texttt{?$\alpha$} a little further. The final event
resolves the metavariable to \texttt{Nat}.

The metavariable is the elaborator's way of carrying the gap between
recorded events. The body \texttt{3} is a recorded event: the user
typed it, and the elaborator has read it. The type \texttt{?$\alpha$} is
unrecorded: the user did not provide it. Between the recorded event
(the body) and the required event (the type), the elaborator carries a
metavariable. The metavariable is constrained but not yet resolved.

This is the algorithmic form of the chapter's question. What can be
said about the gap? The metavariable can be said to be of some sort,
constrained by the surrounding context, but its specific value is not
yet known. The elaborator must carry this state without resolving it
prematurely. If the elaborator resolved \texttt{?$\alpha$} to
\texttt{Nat} as soon as it saw the literal \texttt{3} --- without
considering whether the user might have wanted \texttt{Int} or
\texttt{Real} --- the elaborator would foreclose admissible options.
The correct behavior is to carry the metavariable, accumulate
constraints from the rest of the definition, and resolve at the end
when all constraints are gathered.

The constraints come from many sources. The body provides type
constraints (the literal \texttt{3} demands a numeric type). The
expected return type, if specified, provides another constraint.
Other typeclass constraints (instances that must be available) provide
further restrictions. The default-instance mechanism provides a tie-
breaker when multiple types satisfy all constraints. The elaboration
strategy is to gather all constraints, then resolve.

If the constraints are inconsistent, the elaborator reports an error:
the metavariable cannot be resolved. Example:
\begin{verbatim}
def bar : Int := "hello"
\end{verbatim}
The user has specified the type (\texttt{Int}) and the body
(\texttt{"hello"}). The constraints are: \texttt{"hello"} is a
\texttt{String}; the expected type is \texttt{Int}; \texttt{String}
must equal \texttt{Int}. The unification fails; the elaborator
produces a type mismatch error.

If the constraints are insufficient, the elaborator reports a
metavariable that could not be resolved. Example:
\begin{verbatim}
def baz := []
\end{verbatim}
The body is an empty list. The list has element type \texttt{?$\alpha$},
which is unconstrained. The elaborator cannot infer \texttt{?$\alpha$}
without more information. The error message asks the user to provide an
explicit type annotation.

The metavariable is the elaborator's representation of unresolved
state. It is not an error; it is a normal feature of elaboration. Many
metavariables are carried during the elaboration of a single
definition, and most of them resolve by the end. Only those that remain
unresolved after all constraints are gathered produce errors.

The companion typeclass in \texttt{Episode3.lean} is
\texttt{Metavariable}. The class formalizes the chapter's claim that
metavariables are admissible placeholders. A \texttt{Metavariable}
carries a type (the sort of thing it stands for) and a constraint set
(the conditions it must satisfy when resolved). \texttt{PhysicalProcess}
is the typeclass for processes that produce metavariable outputs ---
outputs that are constrained but not yet resolved. \texttt{COMPARABLE}
certifies that two such processes can be compared: their constraint
sets can be intersected, and the intersection narrows the admissible
resolutions.

Lean's elaborator is doing exactly what the chapter says an
instrument must do. It carries unresolved values through the
computation. It accumulates constraints. It refuses to commit
prematurely. When constraints are sufficient, it resolves. When
constraints are inconsistent, it reports the inconsistency. When
constraints are insufficient, it reports the gap. The elaborator's
behavior is the algorithmic form of admissible measurement under
uncertainty.

This is the chapter's first deep observation. The cubic spline
algorithm carries metavariables (the interior slopes) until the
constraints from the linear system resolve them. The Lean elaborator
carries metavariables (the typed holes) until the constraints from
the body and the context resolve them. The two algorithms are
structurally identical at the level of the chapter's question. Both
carry unresolved values; both resolve from constraint gathering;
both refuse premature commitment.

\section{The Halting Boundary as Interpolation Limit}
\label{se:halting-boundary-as-interpolation-limit}

The cubic spline produced a value at 9:15 AM. The Lean elaborator
resolved \texttt{?$\alpha$ = Nat}. Both algorithms succeeded because their
constraints uniquely determined the metavariable. The next question is
what happens when the constraints do not converge.

Consider a Lean program that should produce a list of natural numbers:
\begin{verbatim}
def search : List Nat := goldbachCounterexamples
\end{verbatim}
where \texttt{goldbachCounterexamples} is a function that searches the
natural numbers for counterexamples to Goldbach's conjecture and returns
the ones it finds. The function is well-defined as a mathematical
specification: for each even integer $n > 2$, check whether $n$ can be
written as a sum of two primes; if not, include $n$ in the list. The
function is also well-typed in Lean: it has type \texttt{List Nat}.

But the function does not terminate. The search is infinite: there are
infinitely many even integers to check, and even if no counterexample
exists, the function must run forever to prove that. Lean's elaborator
cannot evaluate the function during elaboration; the elaboration would
not return.

This is the interpolation limit. The function has a well-defined
mathematical value (a list of natural numbers, possibly empty), but no
algorithmic procedure can produce that value in finite time. The
metavariable representing the function's output cannot be resolved by
the elaborator. The constraints are present and consistent; the
resolution procedure is not.

The \texttt{REPRESENTABLE} typeclass in \texttt{Episode3.lean} marks
this boundary. A function is \texttt{REPRESENTABLE} when there is a
finite algorithmic procedure that produces its value. \texttt{REPRESENTABLE}
typeclass instances are derived for primitive functions (addition,
multiplication, list construction) and propagate through composition
(if $f$ and $g$ are \texttt{REPRESENTABLE}, so is their composition).
The Goldbach search function is not \texttt{REPRESENTABLE}: it
cannot be derived.

What can the elaborator do with a non-\texttt{REPRESENTABLE} function?
Several options. First, accept the function's type definition but not
its evaluation. This is what Lean does by default: the function can be
used in proofs, in type signatures, in the construction of larger
types --- but it cannot be evaluated. The Lean expression
\texttt{\#eval search} would either time out or produce no output.

Second, evaluate the function up to a bound and return a partial
result. This is what \texttt{partial def} permits: the function is
declared partial, indicating that termination is not guaranteed.
Lean trusts the user's assurance and allows the function to be
written. \texttt{\#eval} on a partial function may run forever, but
the elaborator does not refuse to admit the function.

Third, use the function as a metavariable in a larger computation.
The function's value is not resolved; the function is named as an
unknown; downstream computations that depend on it are also
unresolved. The elaborator carries the chain of unresolved values
through the dependent computation.

Each of these is a different protocol for handling the halting
boundary. Each is the chapter's interpolation problem at the algorithmic
limit. The function's value is the metavariable; the procedure to
resolve it does not terminate; the elaborator must decide what to do.
The chapter's lesson is that the elaborator must admit the unresolved
state honestly --- not by silently substituting a default value, not by
silently refusing to admit the function, but by carrying the
unresolved state explicitly and allowing downstream computations to
proceed under the explicit acknowledgment of the gap.

This is the bridge from Chapter 1's halting boundary (the limit of
enumeration) to Chapter 3's interpolation problem (the limit of
metavariable resolution). The two are the same boundary, encountered at
different points in the algorithmic pipeline. Enumeration runs out
when the next mark cannot be produced; interpolation runs out when the
next metavariable cannot be resolved. Both run out at the halting
boundary.

The \texttt{ChaitinsNumberSequence} typeclass in \texttt{Episode3.lean}
formalizes this. The sequence is the binary expansion of $\Omega$, the
halting probability. The $n$-th bit of $\Omega$ encodes the halting
status of program $n$. The bits are well-defined; no algorithm can
compute them all. \texttt{ChaitinsNumberSequence} carries the bits as
metavariables: each bit is named, typed, and bounded (it is 0 or 1),
but its specific value is not resolved.

The compiler can carry the typeclass instance --- the user can write
type signatures involving \texttt{ChaitinsNumberSequence} --- but the
compiler cannot evaluate the specific bits. The bits are the most
precise example of metavariables that are forever unresolved. They
are not unresolved due to insufficient information; they are
unresolved because no algorithm can resolve them. The constraints are
complete; the resolution procedure is undefined.

\texttt{NoisyProcess} extends this to a more general setting. Many
real-world measurements produce outputs that fluctuate around a true
value. The fluctuation is noise; the true value is what the
measurement is intended to characterize. \texttt{PHYSICAL} certifies
that a process produces outputs that include noise. The interpolation
problem for such processes is: from the noisy observations, what is
the best inference about the true value? Linear regression is the
standard answer; the regression's coefficients are metavariables
resolved by least squares.

The noise creates an irreducible gap: even with many measurements,
the regression's coefficients are estimates, not exact values. The
true value is a metavariable that no finite number of observations
can fully resolve --- only narrow. \texttt{COMPARABLE} captures the
property that two such estimates can be compared by the width of
their uncertainty intervals: a tighter interval is comparable to and
strictly less than a looser one.

The chapter's interpolation problem thus has three regimes. In the
favorable regime (cubic spline, Lean type inference for a numeric
literal), the constraints uniquely determine the metavariable, and
the algorithm resolves it cleanly. In the regression regime
(\texttt{PHYSICAL}, noisy data), the constraints narrow the
metavariable to a range, and the algorithm produces an estimate with
named uncertainty. In the limit regime (halting boundary,
\texttt{ChaitinsNumberSequence}), the constraints are well-formed but
no resolution procedure terminates; the metavariable is carried
forever.

All three regimes are admissible in the chapter's sense. The
algorithm must declare which regime it is in, and the calling
context must handle the declared regime. Carrying a Halting-Problem
metavariable as if it were a resolved value is inadmissible. Carrying
an unresolved cubic spline as if the algorithm had failed is
inadmissible. The honest report of the resolution status is part of
the algorithm's output.

\section{Noisy Interpolation}
\label{se:noisy-interpolation}

A scientist measures the same quantity twenty times and obtains twenty
slightly different values. The variations are noise: random fluctuations
due to thermal motion, instrument quantization, electronic interference,
or any of the other small perturbations that affect physical measurement.
The scientist wants the best estimate of the underlying true value.

The standard algorithm is linear regression. Given $n$ data points
$(x_i, y_i)$, the linear regression model assumes $y_i = a + b x_i +
\epsilon_i$, where $\epsilon_i$ is independent Gaussian noise with mean
zero. The parameters $a$ and $b$ are unknown. The procedure finds the
values of $a$ and $b$ that minimize $\sum_i (y_i - a - b x_i)^2$. The
solution is:
\[
b = \frac{\sum (x_i - \bar{x})(y_i - \bar{y})}{\sum (x_i - \bar{x})^2},
\quad
a = \bar{y} - b \bar{x}.
\]
The Python idiom is \texttt{np.polyfit(x, y, 1)}, which returns the pair
$(b, a)$ for a linear fit.

The parameters $a$ and $b$ are the metavariables. Before the regression
runs, they are placeholders constrained by the data. After the
regression, they are resolved to specific values. The values are not
exact: they are the maximum-likelihood estimates under the Gaussian noise
model. The true values of $a$ and $b$ are not measurable; only their
estimates are.

The residual $r_i = y_i - a - b x_i$ is the difference between each
observation and the fitted line. The residual is the irreducible
structure left in the gap after the best admissible linear fit is
removed. If the model is correct, the residuals are random and
uncorrelated; the noise has been characterized and removed. If the
residuals show a systematic pattern --- a curve, a periodic trend, a
heteroskedastic spread --- the linear model is incomplete; a richer
model is needed.

The algorithm's output is the pair $(a, b)$ plus the residual standard
deviation $\sigma$. The standard deviation is the width of the noise
distribution; it characterizes the uncertainty in the fit. A small
$\sigma$ means the data tightly constrains the parameters; a large
$\sigma$ means the data is consistent with a wide range of parameters,
and the estimate is less reliable.

This is the chapter's question for noisy data. The gap between
measurements is filled by the linear fit, but the fit is itself a
metavariable with uncertainty. The uncertainty cannot be eliminated; it
can only be characterized. The algorithm's honest output includes the
uncertainty as part of the answer.

\texttt{NoisyProcess} in \texttt{Episode3.lean} formalizes this. A
process is \texttt{NoisyProcess} when it admits a noise model: a
specification of the distribution of fluctuations around the true value.
\texttt{PHYSICAL} certifies that the process produces outputs that
include noise. A typeclass instance of \texttt{PHYSICAL} for a process
requires the user to provide the noise model. Without the noise model,
the process's outputs cannot be compared meaningfully; with the noise
model, the process's outputs become \texttt{COMPARABLE} (comparable
under uncertainty).

The interpolation problem for \texttt{NoisyProcess} is different from
the interpolation problem for noiseless data. Cubic spline assumed the
data was exact; it forced the spline to pass through every data point.
Linear regression does not force the line to pass through every data
point; the line is the best fit, which generally has nonzero residuals.
The two algorithms differ in their treatment of the data: the spline
believes the data; the regression believes the data plus the noise
model.

Choosing between these two algorithms is a matter of which assumption
is appropriate. If the measurements are highly accurate and the
underlying function is smooth, the cubic spline is appropriate. If the
measurements are noisy and the underlying relationship is approximately
linear, the regression is appropriate. The chapter's lesson is that the
choice of algorithm corresponds to a choice of noise model, and the
noise model is part of the calibration of the instrument.

The compiler analog of noisy interpolation is the parsing of
ambiguous source code. When the parser encounters an ambiguity (e.g.,
operator precedence that could be interpreted multiple ways), it
must select one interpretation. The selection is a function of the
parser's grammar plus the parser's tie-breaking rules. The
interpretation is a metavariable resolved by the parser's heuristics;
the heuristics are the parser's calibration. Two different parsers,
with different calibrations, may resolve the same ambiguous expression
differently.

In well-designed languages (Lean among them), the grammar is
unambiguous, so the parser does not need heuristics: every valid
source produces a unique parse tree. In less carefully designed
languages, the heuristics produce surprises. The chapter's point is
that the parser's behavior on ambiguous input is the algorithmic
analog of the regression's behavior on noisy data: when the input
does not uniquely determine the output, the algorithm must declare
its tie-breaking rule, and the output must be interpreted relative
to that rule.

\section{ChaitinsNumberSequence in Lean}
\label{se:chaitinsnumbersequence-in-lean}

The Lean file \texttt{Episode3.lean} defines a type whose existence the
compiler can confirm but whose specific values the compiler cannot
compute. This is \texttt{ChaitinsNumberSequence}. The definition has the
shape:

\begin{leanexcerpt}{ChaitinsNumberSequence}
inductive ChaitinsNumberSequence
  | halting    : Fact $\to$ Computation $\to$ ChaitinsNumberSequence
  | nonhalting : Fact $\to$ Computation $\to$ Computation
               $\to$ ChaitinsNumberSequence
\end{leanexcerpt}

The inductive has two constructors. The \texttt{halting} constructor
records a \texttt{Fact} and a single \texttt{Computation}: the case
where the program halts (and the bit is admissibly $1$). The
\texttt{nonhalting} constructor records a \texttt{Fact} and two
\texttt{Computation}s: the case where the program does not halt under
either of two candidate witnesses (and the bit is admissibly $0$). The
inductive together specifies the bit-by-bit construction of Chaitin's
constant $\Omega = \sum_{p \text{ halts}} 2^{-|p|}$, with the halting
and nonhalting layers stacked as the sequence is read.

The compiler can elaborate this type definition. Lean's type checker
verifies that the type is well-formed: \texttt{bits} is a function of
the correct type; \texttt{halting\_property} is a predicate that
correctly relates \texttt{bits} to the halting status of programs. The
type elaborates. It is a valid Lean type.

The compiler cannot construct an instance. To produce a value of type
\texttt{ChaitinsNumberSequence}, the user would need to provide a
specific \texttt{bits} function plus a proof that it satisfies the
halting property. The proof requires deciding, for every $n$, whether
program $n$ halts. The decision is the Halting Problem. No algorithm can
provide such a proof for all $n$. The constructor is not callable.

The type exists; values of the type do not. This is the cleanest
algorithmic expression of the chapter's interpolation limit. The
metavariable (the \texttt{bits} function) is named, typed, and
constrained. Its value is not resolvable by any finite procedure. The
elaborator can carry the type through the namespace, the kernel, and
all type-level operations. The elaborator cannot evaluate any specific
bit.

This may seem like a paradox. How can a type exist without values? In
Lean, the existence of a type is a syntactic question: the type is
well-formed if its definition type-checks. The existence of a value is
a semantic question: a value exists if it can be constructed. The two
questions are distinct. A type may type-check without admitting any
constructible value.

The simplest example is the empty type \texttt{Empty}: well-formed,
type-checks, but has no values (no \texttt{Empty} can be constructed).
\texttt{Empty} is useful: it serves as a witness of impossibility (a
proposition is false if it implies $\text{Empty}$). The type
\texttt{ChaitinsNumberSequence} is similar: it is well-formed, it
type-checks, but no constructive procedure can produce a value. Its
usefulness is also as a witness: it names the boundary of computable
sequences.

The compiler's response to \texttt{\#check
ChaitinsNumberSequence} is a confirmation that the type exists:
\texttt{ChaitinsNumberSequence : Type}. The response to \texttt{\#eval
(ChaitinsNumberSequence.mk ...)} would be an error: the constructor
cannot be called without providing the halting-property proof.

The book's argument turns on this distinction. The Lean compiler can
admit types whose values are not effectively constructible. The
admission is the type-check, which is a finite procedure. The
construction is the evaluation, which is not a finite procedure. The
book uses this distinction to argue that the universe, which must
admit types of the same kind, is doing something the compiler cannot
simulate: the universe constructs values of types like
\texttt{ChaitinsNumberSequence}, while the compiler can only check
that the type is well-formed.

This is the conclusion the book is building toward, stated for the
first time here in algorithmic form. The selection rule that picks
one admissible continuation requires resolving a metavariable whose
constraints are well-formed but whose resolution procedure does not
halt. The metavariable is a sequence of bits like
\texttt{ChaitinsNumberSequence}. The constraints are the
admissibility conditions of the continuation. The resolution
procedure --- the algorithm that picks which bit is 0 and which is 1
--- would be a Halting-Problem solver if it existed. It does not.

In Lean, the user can prove theorems about
\texttt{ChaitinsNumberSequence} without ever constructing an instance.
Theorems about the type's properties (e.g., that $\Omega$ is
algorithmically random) are provable from the definition; theorems
about specific bit values are not. The type-level reasoning is
admissible; the value-level computation is not.

This is the chapter's closing observation. The compiler can carry
metavariables, including metavariables whose specific values are
provably unresolvable. The carriage is admissible; the resolution may
not be. The honest report of the algorithm's output therefore must
distinguish between two kinds of metavariables: those that the
algorithm resolves (cubic spline slopes, Lean type inference for a
numeric literal), and those that the algorithm carries without
resolving (Halting-Problem instances,
\texttt{ChaitinsNumberSequence} bits). Both are admissible. They are
not the same.

\begin{bridge}

The chapter's question was what can be said about the gap between
recorded events without inventing structure not in the ledger.

The answer is: the gap is filled by metavariables, and the metavariables
are constrained by the surrounding record. Cubic spline fitting resolves
the metavariables to specific values through a tridiagonal linear
system. Lean's elaborator resolves the metavariables to specific types
through unification. Noisy interpolation resolves the metavariables to
estimates with named uncertainty. The halting boundary names
metavariables whose resolution procedure does not terminate; the
\texttt{ChaitinsNumberSequence} type carries such metavariables as
named-but-unresolved entities. Each algorithm declares its resolution
status as part of its output.

What remains is that single records, with their gaps filled, do not yet
combine into multi-record structures. Two records may be about different
aspects of the same physical system, or about the same system at
different times, or about systems whose physical relationship is part
of the inquiry. Chapter 4 asks how multiple records combine into a
joint record without collapsing into one undifferentiated heap.

\end{bridge}

\begin{coda}{The Legal Contract with Blanks}

The lawyer's office overlooks the harbor. Two parties are seated at the
table: the buyer's representative, the seller's representative, the
notary, and the lawyer. A document has been printed and bound: thirty-two
pages, double-spaced, in the standardized template the lawyer's firm
uses for commercial real-estate transactions.

The document is complete except for three blanks.

Blank one is on the cover page: the effective date of the contract.
Today is October 24th, but the contract will not become effective until
all three signatories have signed. The effective date will be the date
of the last signature.

Blank two is on page eleven, paragraph fourteen: the buyer's signature.
The buyer has reviewed the document and has agreed to its terms, but
the lawyer has identified a tax provision in paragraph nine that needs
clarification before the buyer commits. The buyer is willing to sign,
but not until paragraph nine is resolved.

Blank three is on page fourteen, paragraph eighteen: a placeholder for
the seller's initials. The seller has read the document twice and has
agreed to all terms. The initial is a formality. The seller will sign
once the buyer signs.

The lawyer walks through the document line by line. The structure is
admissible: the clauses are well-formed, the obligations are clearly
specified, the remedies are enforceable. The contract is a
\texttt{REPRESENTABLE} legal instrument. It has type \texttt{Contract}.
It would type-check.

The three blanks are metavariables.

The effective-date blank is constrained: it must be a date in the
future, no earlier than today, no later than thirty days from today
(after which the offer expires). It is also constrained by the structure
of the contract: it must be a date on which all signatories signed. The
metavariable's resolution depends on the resolution of the other two
metavariables (the signatures), and the resolution is automatic ---
the moment the second signature is added, the effective date is the
date of the second signature.

The buyer's signature blank is constrained: it must be the buyer's
authentic signature, executed in the lawyer's presence, witnessed by the
notary. The resolution depends on the resolution of the tax
clarification, which is a separate metavariable in the system but not
in the document. The buyer's lawyer is researching the tax provision;
the result will be communicated by phone within the hour. Until then,
the buyer's signature blank is carried as unresolved.

The seller's initial blank is constrained: it must be the seller's
authentic initial, executed by the same procedure. The resolution
depends on the buyer's signature: the seller has stated they will
initial once the buyer signs. The seller's initial is a deferred
resolution; it is not unresolved due to insufficient information, but
due to a chosen sequencing.

This is the algorithmic structure of Chapter 3 in legal form. The
contract is a \texttt{REPRESENTABLE} instrument with three
metavariables. The metavariables are constrained. They are not yet
resolved. They are not random; they are not arbitrary. They have type
(date, signature, initial) and admissibility conditions (date in
window, witnessed signature, witnessed initial).

The lawyer's role is to carry the metavariables admissibly. She does
not pretend the contract is complete. She does not allow the buyer to
sign before the tax provision is resolved. She does not pretend the
seller has already initialed when the seller has not. The metavariables
are present in the record exactly as such: blanks with named
constraints. The contract carries them as gaps, not as values.

The lawyer makes a phone call to the buyer's tax counsel. The
counsel returns the call eighteen minutes later. The tax provision in
paragraph nine, when read in conjunction with subsection 5.2.1 of the
state's revenue code, requires a small modification: a clarifying
clause specifying that capital gains tax will be paid by the seller
out of the proceeds before transfer to the buyer. The lawyer drafts
the amendment in fountain pen on a separate page, has it initialed by
both parties as a contract addendum, and re-sequences the affected
sections.

The buyer reviews the amendment. The amendment resolves the
constraint. The buyer signs paragraph fourteen. The signature is
witnessed by the notary. Blank two is now resolved.

The seller, seeing the buyer's signature, initials paragraph eighteen.
The initial is witnessed by the notary. Blank three is resolved.

The notary records the date and time of the second witnessed event:
October 24th, 3:47 PM. This is the effective date. The lawyer writes
the date on the cover page. Blank one is resolved.

The contract is now complete. All three metavariables have been
resolved. The contract has transitioned from \texttt{REPRESENTABLE}
(well-formed but unsigned) to \texttt{EXECUTED} (signed and
binding). The transition is not a measurement of anything physical;
it is the resolution of three placeholders into three specific
values.

The lawyer carries a second metavariable that does not appear in the
document: her assessment of the buyer's understanding. The buyer is
sophisticated, but the tax provision required research. The lawyer
believes the buyer understands the contract; she has no proof. The
assessment is carried as a \texttt{Metavariable} in the lawyer's own
record-keeping. If a dispute arises later (the buyer claims the tax
provision was unclear), the assessment will be retrieved from the
lawyer's notes, along with the timestamp of the buyer's
consultation with their tax counsel. The lawyer's notes are the
audit trail. The metavariable is constrained by the audit trail, and
its resolution is what the trail says.

If the buyer had not signed --- if the tax clarification had not
resolved --- the contract would have expired in thirty days. The
expiry is the halting boundary: the metavariable could not be
resolved in the available time, so the contract is closed in the
unresolved state. The unsigned contract is not erased; it is filed
as a non-execution. The record of the attempt remains. The buyer and
seller may re-engage in the future under a new contract.

This is the chapter's point in legal form. The metavariable is the
admissible representation of unresolved state. The contract carries
it without pretending it is resolved. The resolution proceeds by
gathering constraints (the tax clarification, the buyer's review,
the seller's commitment). The resolution may complete (signatures
obtained) or may halt unresolved (offer expires). Either way, the
record is honest about what was resolved and what was not.

The notary closes the day's register. October 24th has produced
seven witnessed acts. The day's record is complete. Tomorrow's
register will have its own metavariables, its own contracts, its own
gaps. The notary's record is the \texttt{Study}: the structured
accumulation of witnessed events. Each event has admissible
metadata; each metavariable has an admissible resolution status. The
record grows by accumulation; it does not lose any prior event.

\end{coda}
